% Generated draft: requires mathematical and source review.
\documentclass[11pt]{article}
\usepackage[margin=25mm]{geometry}
\usepackage{fontspec}
\setmainfont{Latin Modern Roman}
\usepackage{amsmath,amssymb,mathtools,unicode-math}
\setmathfont{Latin Modern Math}
\usepackage{enumitem,needspace,xurl,hyperref}
\setlength{\parindent}{0pt}
\setlength{\parskip}{5pt}
\setlength{\emergencystretch}{3em}
\newcommand{\N}{\mathbb{N}}
\newcommand{\Z}{\mathbb{Z}}
\newcommand{\Q}{\mathbb{Q}}
\newcommand{\R}{\mathbb{R}}
\newcommand{\C}{\mathbb{C}}
\newcommand{\F}{\mathbb{F}}
\newcommand{\A}{\mathbb{A}}
\newcommand{\PP}{\mathbb{P}}
\newcommand{\E}{\mathbb{E}}
\newcommand{\eps}{\varepsilon}
\newcommand{\dd}{\,\mathrm{d}}
\begin{document}
\section*{UnsolvedMath ID: 30001034}
\textbf{Source problem code:} catalogue entry
\subsection{Definitions and mathematical statement}
Let $f:\Z_{\geq1}\to\Z_{\geq1}$ be nondecreasing and satisfy $f(n)\geq\exp(n)$ for every $n\geq1$. For a finitely generated group $G$ and a finite symmetric generating set $S$, define the outer boundary $\partial_S A=SA\setminus A$ of a finite subset $A\subseteq G$, and use the convention \[\operatorname{Fol}_{G,S}(n)=\min\{|A|:0<|A|<\infty,\ |\partial_S A|\leq |A|/n\}.\] Set the minimum to $+\infty$ if no admissible set exists. Having finite values for all $n$ is the F\o lner amenability condition in this convention. The source asks whether every such $f$ is realized by some finitely generated group: \[\operatorname{Fol}_G(n)\asymp f(n).\] No smoothness, doubling, regularity, or computability condition is imposed on the prescribed function beyond monotonicity and the exponential lower bound. Any realizing group would have finite F\o lner values and thus be amenable; a finite group could not realize the required unbounded growth. The source suppresses the generating set and does not define $\asymp$. Its usual pointwise comparison meaning would be $cf(n)\leq\operatorname{Fol}_{G,S}(n)\leq Cf(n)$ for all sufficiently large $n$ and constants $c,C>0$. In group geometry one also uses the weaker comparison permitting rescaling of the argument, for example $F(n)\leq C f(\lceil Cn\rceil)$ and the reverse inequality with the roles exchanged. These versions need not agree for arbitrary rapidly growing or irregular $f$, and the latter is the usual setting for independence of generating-set choices. The requested realization cannot be assigned one of these meanings silently. The boundary convention is similarly explicit in this draft, not claimed as the report's only convention. The supplied 2026-08-21 triage labels the question open and reports broad constructions and characterizations with extra regularity, without a realization theorem for every function in the stated class.
\subsection{Short English statement}
Can every nondecreasing function bounded below by $e^n$ occur as a finitely generated group's F\o lner function? The comparison relation and generating-set/boundary conventions matter, especially because no regularity of the prescribed function is assumed.
\subsection{Sources}
\begin{itemize}\item Local UnsolvedMath registry, permanent ID 30001034, source problem catalogue entry, \emph{Realization of Rapidly Growing F\o lner Functions}, proposed year 2008; no individual proposer supplied.\item Supplied citation: \emph{Geometric Group Theory, Hyperbolic Dynamics and Symplectic Geometry} (2009), Oberwolfach Reports, DOI 10.4171/owr/2008/42, pp.\ 2337--2406: \url{https://doi.org/10.4171/owr/2008/42}. Supplied literature label: Oberwolfach report.\item Supplied literature label: Erschler--Zheng, Folner functions, \url{https://arxiv.org/abs/1708.04730}.\item Supplied literature label: Cavaleri, prescribed growth, \url{https://arxiv.org/abs/1703.04133}.\item Literature assessment date: 2026-08-21. The status and progress claims are supplied dated curation, not independently verified here.\end{itemize}
\end{document}
